\vspace{0.5em}\noindent\textbf{TỪ ZERO ĐẾN ZERO-DOWNTIME: DEVOPS \& CI/CD TRÊN
AWS}\par

Khi bắt đầu DevOps trên AWS, mình từng nghĩ CI/CD chỉ là viết script tự
động đẩy code lên máy chủ. Nhưng các hệ thống Production thực tế đã dạy
rằng DevOps thực chất là về \textbf{Độ tin cậy}, \textbf{Bảo mật} và
\textbf{Tốc độ}.

\vspace{0.5em}\noindent\textbf{3 Bài Học Quan Trọng}\par

\vspace{0.5em}\noindent\textbf{1. Infrastructure as Code (IaC): Nói Không Với
ClickOps}\par

\begin{itemize}
\tightlist
\item
  Bấm tay trên console lúc đầu nhanh nhưng là ác mộng khi tái lập môi
  trường hoặc audit lỗi.
\item
  \textbf{Dùng Terraform hoặc AWS CDK}. Coi hạ tầng như code. Mọi thay
  đổi qua Pull Request \(\rightarrow\) Plan \(\rightarrow\) Code Review.
\end{itemize}

\vspace{0.5em}\noindent\textbf{2. Shift-Left Security Trong CI/CD
Pipeline}\par

\begin{itemize}
\tightlist
\item
  Đừng đợi lên Production mới quét lỗ hổng.
\item
  \textbf{Tích hợp bảo mật sớm:}

  \begin{itemize}
  \tightlist
  \item
    Quét file IaC với Checkov/TFLint
  \item
    Quét container image với Amazon ECR Image Scanning
  \item
    Quản lý secrets tập trung với AWS Secrets Manager (không hardcode)
  \end{itemize}
\end{itemize}

\vspace{0.5em}\noindent\textbf{3. Blue/Green \& Canary Deployments Cho
Zero-Downtime}\par

\begin{itemize}
\tightlist
\item
  Deploy kiểu in-place dễ gây gián đoạn và lỗi 502.
\item
  \textbf{Dùng AWS CodeDeploy:}

  \begin{itemize}
  \tightlist
  \item
    \textbf{Blue/Green}: Dựng môi trường song song, test rồi chuyển
    traffic
  \item
    \textbf{Canary}: Chuyển 10\% traffic sang bản mới, theo dõi
    CloudWatch Metrics rồi rollout 100\%
  \end{itemize}
\end{itemize}

\vspace{0.5em}\noindent\textbf{Kiến Trúc Pipeline Chi
Tiết}\par

Dưới đây là luồng CI/CD đã được áp dụng trong thực tế:

\textbf{Bước 1: Developer commit code lên GitHub}

Khi lập trình viên đẩy code lên repository, đây là điểm khởi đầu của
toàn bộ pipeline.

\textbf{Bước 2: AWS CodePipeline phát hiện thay đổi và kích hoạt
pipeline}

CodePipeline liên tục theo dõi repository. Khi có commit mới, pipeline
tự động được kích hoạt mà không cần can thiệp thủ công.

\textbf{Bước 3: AWS CodeBuild chạy Unit Tests}

Kiểm tra chất lượng code bằng các bài test tự động. Nếu test thất bại,
pipeline dừng lại ngay để tránh deploy code lỗi.

\textbf{Bước 4: CodeBuild chạy Security Scan}

Quét mã nguồn và các dependency để phát hiện lỗ hổng bảo mật. Phát hiện
sớm giúp giảm chi phí sửa lỗi so với khi đã lên Production.

\textbf{Bước 5: Build container image và push lên Amazon ECR}

Tạo Docker image từ source code, sau đó đẩy lên Amazon Elastic Container
Registry để lưu trữ và quản lý các phiên bản image.

\textbf{Bước 6: ECR tự động quét image}

Ngay khi image được push lên, ECR tự động quét để phát hiện CVE trong
các layer của image. Kết quả quét được ghi lại để đánh giá trước khi
triển khai.

\textbf{Bước 7: CodeDeploy (hoặc Helm) triển khai lên Amazon EKS}

Image được triển khai lên Kubernetes cluster. Tùy vào cách thiết lập, có
thể dùng CodeDeploy tích hợp sẵn với EKS hoặc dùng Helm để quản lý ứng
dụng dạng charts.

\textbf{Bước 8: Áp dụng chiến lược Blue/Green hoặc Canary}

Thay vì deploy trực tiếp lên production, pipeline sử dụng một trong hai
chiến lược:

\begin{itemize}
\tightlist
\item
  Blue/Green: Tạo bản Green (mới) song song với bản Blue (hiện tại). Sau
  khi test kỹ trên bản Green, traffic được chuyển từ Blue sang Green.
\item
  Canary: Chuyển một phần nhỏ traffic (ví dụ 10\%) sang bản mới, theo
  dõi các chỉ số trước khi chuyển toàn bộ.
\end{itemize}

\textbf{Bước 9: Theo dõi metrics trên CloudWatch và tự động rollback}

Trong suốt quá trình deploy, CloudWatch theo dõi các chỉ số quan trọng
như tỷ lệ lỗi, thời gian phản hồi. Nếu phát hiện bất thường, pipeline tự
động rollback về bản cũ để giảm thiểu tác động.

\vspace{0.5em}\noindent\textbf{Giải thích từng thành
phần:}\par

\textbf{GitHub:} Nơi lưu trữ source code, tích hợp sẵn với CodePipeline.

\textbf{AWS CodePipeline:} Dịch vụ orchestration CI/CD, điều phối toàn
bộ luồng từ build đến deploy.

\textbf{AWS CodeBuild:} Dịch vụ build serverless, thực hiện test,
security scan và build image.

\textbf{Amazon ECR:} Registry container image, tích hợp quét CVE tự
động.

\textbf{AWS CodeDeploy / Helm:} Công cụ triển khai với các chiến lược
Zero-Downtime.

\textbf{Amazon EKS:} Kubernetes cluster nơi ứng dụng chạy.

\vspace{0.5em}\noindent\textbf{Kết quả đạt
được:}\par

\begin{itemize}
\tightlist
\item
  Thời gian deploy giảm từ 30 phút xuống dưới 4 phút
\item
  Phát hiện sớm lỗ hổng bảo mật ngay từ khâu build
\item
  Triển khai không downtime nhờ Blue/Green và Canary
\item
  Tự động rollback nếu có vấn đề sau deploy
\end{itemize}

\vspace{0.5em}\noindent\textbf{Điểm Quan Trọng}\par

Công cụ chỉ chiếm 20\%. 80\% còn lại là \textbf{Thiết kế Pipeline} và
văn hóa \textbf{Release nhỏ, liên tục}.

\textbf{Bài viết gốc:}
\href{https://www.facebook.com/groups/awsstudygroupfcj/permalink/2229176364513990/?rdid=aT7nNGRCmytCuj4K\#}{Facebook
AWS Study Group FCJ}
